\section*{Scope and research status}
\addcontentsline{toc}{section}{Scope and research status}
This paper treats the supplied Alp\"oge and OpenAI theorems as cited
results, as requested. It distinguishes those premises from deductions
proved in this run. ``Finding'' means such a deduction, including a
counterexample or an elementary sharpened form of a known result; it does
not assert historical priority. Exact source locators accompany the
statements used. The accompanying check register describes what was
actually checked and does not promote a text scan or a numerical test to
a formal proof certificate.

The main negative conclusions are P-F31 and P-F57. Both clauses of
ordinary plethysm PH3 in GCT VI, Hypothesis~1.6.7, are false;
the narrower polynomial positivity-index and strict-saturation-index
bounds in Hypotheses~1.6.6 and 1.6.5(a) are false as well.
For the actual cubic family, P-F57 proves
$p(f_n)>2^{cn}$ and $s(f_n)\ge2^{cn}$ for every fixed
$0<c<1/4$ and all sufficiently large $n$.
An explicit version gives these bounds with $c=1/8$ for
every $n\ge2\cdot10^{80}$. The source parameter is
$h=\binom{n+2}{3}^2+n+1=\Theta(n^6)$.
These are the literal constituent-polynomial indices; the
separate density clause SH(b), Kronecker-specific polynomial
indices, and the complexity-theoretic decision hypotheses
are not refuted by this result. The family has
the single inner shape $(3)$. Its invariant Hilbert series has an
actual simple pole at every primitive $2^{n-1}$st root, although the
source's parameter is $O(n^6)$. Every rational expression with
polynomial numerator and factors $1-t^a$ in its denominator must
therefore use some $a\ge2^{n-1}$. P-F30 supplies the complete
support classification that excludes every competing nonzero fixed
component; P-F23, P-F24 and P-F28 then prove the pole and its leading
coefficient. P-F19 had already refuted the canonical parameter clause.
The pole theorem alone does not give the index refutation.
P-F56 isolates its primitive dyadic wave, and P-F57 supplies
the additional global comparison that does.
P-F62 preserves the exponential strict-index obstruction after
specified polynomial odd dilations.  The resulting actual inputs
have positive initial multiplicity and a positive value at every
positive stretch, yet still have exponential literal indices.
This strengthens the distinction between actual nonvanishing and
the off-residue constituent tests; P-F48's density limitation remains.

There is now a separate Kronecker conclusion.  P-F60 constructs a
balanced chain tensor with exact product-determinant stabilizer
$\boldsymbol\mu_{2^{n-1}}$, and identifies its invariant ring,
including multiplication, with the actual canonical algebra for
$k((s^n),(s^n),(s^n))$.  Every homogeneous system of parameters
contains a degree divisible by $2^{n-1}$, whereas the unified source
rank is $n^2+5n$.  Thus the canonical clause of GCT VI
Hypothesis~3.3.4 is false for these triple-rectangular Kronecker
inputs.  The numerical positive-form clause and strict Kronecker
indices remain open: the constructed smooth fixed curve has an
explicit local conormal factor, but other global fixed components
have not been controlled. P-F64 proves that a second positive-dimensional
fixed component exists, refuting the unique-component shortcut.
P-F65 computes its full stabilizer and local determinant and
proves that these two explicit components have the same phase at
the principal dyadic root.  Their mutual cancellation is excluded;
the sum over all components remains open.
P-F58 excludes specified direct
plethysm transfers; P-F59 gives another actual Kronecker family
with efficiently constructible positive forms and polynomial indices
despite exponential minimum periods.

A central positive result obtained here is P-F11: a global
quasi-polynomial with nonnegative rational Hilbert numerator and
$r$ denominator factors of degrees at most $A$ has positivity index
less than $rA+128r^3A^3$. P-F12 applies this to show that the revised
PH3 of GCT VI implies its revised PH2. This is an implication from
bounded positive Hilbert data, not a proof that the requested
denominator bounds exist.
P-F14--P-F15 separately refute the rank-only PH3, PH2 and saturation
bounds in the broader subgroup-restriction formulation: faithful
polynomial torus embeddings have fixed group ranks and unbounded
indices as their character weights vary. The narrower plethysm
formulation is not refuted by those examples. P-F63 separately
refutes the broad subgroup density clause under finite
constant-additive source specifications, by fixed product padding.
The narrow ordinary-plethysm density assertion remains open.
P-F67 gives a sufficient amplification criterion inside its exact
raw-part bitlength population: logarithmic length overhead and
polynomial fibers would already refute that density assertion.
Such a padding map has not been constructed.
P-F16--P-F18 prove finite uniform bounds for a fixed semisimple
embedding and, for ordinary $GL_n$ plethysm, for fixed inner dimension.
P-F19 then gives an exponential lower bound for that optimal canonical
modular-index envelope, using the single inner shape $(3)$.
It refutes the stronger canonical clause of ordinary plethysm PH3.
P-F31 additionally refutes the assertion about some positive rational
Hilbert expression. P-F57 subsequently resolves the narrower
ordinary plethysm index questions negatively by a separate argument.
P-F20 distinguishes a large universal denominator from an ordinary
Hilbert-series pole. P-F21--P-F22 refute general transfers from large
parameter degrees even for rational Gorenstein domains with one fixed
Hilbert series. P-F23 supplies a precise noncancellation criterion;
P-F24 proves that the chain gives a smooth one-dimensional fixed
component in the actual cubic quotient. P-F26 confines every competing component
to the discriminant boundary. P-F27 proves that any such component
requires three-distinct-variable monomials in every finite-lift
eigenbasis; even a connected graph part forces the smooth chain.
P-F29 excludes the complementary cube-and-squarefree support class,
so any remaining supports would have to mix the two types with
disconnected graph part. P-F30 closes that final support case,
including rank-deficient supports and arbitrary two-power lift order.
P-F28 gives the exact local chain coefficient and its cyclotomic
valuation; P-F31 identifies it with the actual leading coefficient.
P-F32 supplies a separate exterior-duality reduction and an exact
first nonzero stretch on an infinite cubic family, giving a
$\Theta(\sqrt n)$ lower bound compatible with polynomial PH2 and SH.
P-F33 prevents an inference from large period and Gorenstein
reciprocity alone to large indices. P-F34 proves a large Fourier
amplitude for the actual cubic family and reduces one route to
an exponential index lower bound to a specified, still unproved
bound on the residue-average polynomial. P-F35 computes both
indices exactly as $23$ in the four-variable case.
P-F56 proves that the maximal two-part of a Hilbert pole is exactly
$2^{n-1}$ and isolates those simple poles by a half-period difference
of constituents. A single positive wave value exceeding the actual
coefficient gives a strict-index lower bound. A specified pair has
quartic exponential amplitude; an explicit bound on the total remaining
conjugates yields exponential indices. P-F57 proves that bound.
The part indexed by normalized multipliers at most $2^{\delta n}$
is proved negligible for every fixed $\delta<1/600$.
A contraction uses only quadratic-many factors at each dimension
increment. Its near-multiplier estimate and far-multiplier
dyadic counting estimate prove the uniform inequality in P-F57.
Positive Hilbert forms also yield
shifted polytope models (P-F2), and the main counterexample ray admits
an optimal two-dimensional model (P-F1).

The converse direction from qualitative purity to the required counting
and complexity properties fails in concrete examples. PLE-F7 supplies a
normal surface with pure even Tate ordinary and intersection cohomology
whose Hilbert function violates homogeneous PH1. PLE-F8 identifies it
with the actual canonical plethysm multiplicity algebra on that ray,
including multiplication, and supplies a basis with nonnegative
multiplication constants. P-F10 keeps the smooth
underlying variety fixed while its grading indices grow arbitrarily.
P-F7 shows that finite smooth proper pure Tate schemes can encode all
\#3SAT counts when specified by equations. These examples identify
missing hypotheses; they do not prove a logical independence theorem
between RH and any GCT assertion.

The exact theta input also gives a concrete finite-operation
obstruction. P-F68 proves polynomial coefficient growth for fixed
finite definite-theta formulas and ordinary finite lambda operations,
excluding the actual input $\Theta_{E_6}/\delta$ and quantifying the
rank or polynomial-degree growth needed for varying formulas.
In the nonstandard cubic example, P-F66 excludes the split source
numbers as geometric Frobenius eigenvalues. P-F69 supplies a
different exact geometric comparison: at each $q=3^e$, they are
rationally normalized real-quadratic traces on the first cohomology
of a simple ordinary abelian surface over $\mathbb F_{q^{116}}$.
This scalar fiberwise comparison does not identify the source
products or the supplied theta operations.
P-F70 supplies an additional actual source-block map: both
specialized generators and every ordered word act by rational
quasi-endomorphisms on pure $H^1$, with the prescribed central
action Rosati positive.  P-F72 uses a different construction to
make each source number an exact component Frobenius trace,
with an ordinary Tate twist and a pure weight-$65$ projected
cohomology space.  The construction has $K$-rank $2^{117}$
and fails base-extension compatibility.  Formal graded
canonical coefficient objects and the supplied theta comparison
remain open.

No separation or equality of $\mathsf P$ and $\mathsf{NP}$, or of the
usual complex algebraic classes $\mathsf{VP}$ and $\mathsf{VNP}$, is
proved. The remaining research targets are displayed separately in
Part~II and in \texttt{QUESTIONS.tsv}. A statement is not marked proved
merely because its missing input has been named.

References to numbered GCT VI hypotheses concern the supplied
arXiv:0704.0229v4 text. The broad subgroup-restriction formulation
of Problem 1.1.3 and Sections 3.3--3.4 is distinguished throughout
from irreducible-inner Problem 1.1.2 and its ordinary $GL_n$
plethysm subfamily. Introductory notational inconsistencies about
the canonical invariant algebra are resolved using the explicit
construction and modular-index definition in Sections 4.1.1--4.1.2.

\subsection*{The supplied analytic statements}
The five theorem declarations in \texttt{ALPOGE\_Main.lean} are the
precise analytic premises. Their types are imported from the supplied
\texttt{ALPOGE\_Statements.lean}, except for Mathlib's
\texttt{RiemannHypothesis} itself.
\begin{longtable}{@{}p{.48\linewidth}p{.46\linewidth}@{}}
\toprule
Theorem name (statement identifier) & Scope used here\\
\midrule
\endhead
\path{riemannHypothesis} (A-RH) & Mathlib's named Riemann-hypothesis proposition.\\[3pt]
\path{riemannZeta_eq_zero_re} (A-ZETA) & A zero of $\zeta(s)$, other than a negative even integer and with $s\ne1$, has $\Re s=1/2$.\\[3pt]
\path{LFunction_ne_zero_of_half_lt_re} (A-DIR-NV) & For every positive modulus and Dirichlet character, $L(\chi,s)\ne0$ if $\Re s>1/2$ and $(\chi\ne1\text{ or }s\ne1)$.\\[3pt]
\path{LFunction_eq_zero_re} (A-DIR-Z) & All Dirichlet $L$-zeros in $0<\Re s<1$ have real part $1/2$.\\[3pt]
\path{LFunction_eq_zero_re_of_isPrimitive} (A-DIR-PRIM) & The corresponding statement restricted to primitive Dirichlet characters.\\
\bottomrule
\end{longtable}
These declarations quantify over Dirichlet characters, not arbitrary
number fields or their Hecke characters. C-GRH-01 proves the exact
cyclotomic-field extension needed below. The present working set contains
the theorem wrappers, their statement file and a supplied scan report;
it does not contain all imported proof modules. No Lean kernel rebuild
was performed. This records the scope of this run's verification, without
reopening the requested treatment of the theorems as cited results.

\subsection*{The positivity method in the supplied note}
The determinant sign $W(1_S)\ge0$ is in \emph{Remark}, p.~3, numbered
item~4. The sign-only discussion of finite-field purity is in footnote~3
on that page; numbered item~5 concerns switching to the von Mangoldt
sum. GGP periods and interlacing appear in item~7, iterated theta lifts in
item~8, and the Bessel-period seesaw in item~12 on p.~4. We use these
locators rather than conflating the footnote with numbered item~5.

The transfer question is mathematical: identify the representation
multiplicity or graded coordinate-ring quotient with the object to which
a sign or purity theorem applies, then supply quantitative control.
P-F3 constructs a positive determinant for each coefficient
unconditionally, exposing the dimension cost. P-F4 separates pointwise
sign from character positivity. P-F5 distinguishes standard interlacing
from restriction along a Schur-functor homomorphism. P-F6 states the
precise categorical data sufficient for positivity of a based product.
Neither the four-page note nor the five theorem types supplies those GCT
identifications automatically.
P-F36 supplies a more specific requirement in the actual nonstandard
cubic example: a quadratic coefficient extension survives every
finite power regrading. Rational projectors cannot isolate its
conjugate blocks, and finite graded Laurent characters cannot
produce the displayed raw coefficient. A proposed transfer must
account for this descent datum; the relaxed conjectures themselves
are not refuted.

\subsection*{Reading the package}
Part~I contains the checked findings and the exact cited background on
which they rely. Part~II contains unresolved targets and unsuccessful
transfer attempts; it asserts no unchecked theorem. The source archive
is identified by a SHA-256 manifest. The audit register records model,
effort disclosure, agent/run, date, method and outcome for each check.
The finite computations are sanity checks, separately labelled from
whole-proof readings and independent re-proofs.
